% ============================================================
% Übungsblatt 1: Enter Numerical Methods
% Zum Vorlesungsskript "Numerical Methods" (Prof. Mark Schutera),
% Kapitel 1 (S. 9-16)
% ============================================================
\documentclass[11pt,a4paper]{scrartcl}

\usepackage[T1]{fontenc}
\usepackage[utf8]{inputenc}
\usepackage[ngerman]{babel}
\usepackage{lmodern}
\usepackage{amsmath}
\usepackage{amssymb}
\usepackage[a4paper,margin=2.5cm]{geometry}
\usepackage{enumitem}
\usepackage{booktabs}
\usepackage[hidelinks]{hyperref}
\usepackage{catppuccin-dark}

% Aufgaben-Zähler, damit \label/\ref trotz \section* funktionieren
\newcounter{aufgabe}
\newcommand{\aufgabe}[2]{%
  \refstepcounter{aufgabe}%
  \paragraph{Aufgabe \theaufgabe: #1 \normalfont\mdseries[#2]}}

\title{Übungsblatt 1: Enter Numerical Methods}
\subtitle{Zum Vorlesungsskript \glqq Numerical Methods\grqq{} (Prof.\ Dr.-Ing.\ Mark Schutera), Kapitel~1, S.~9--16}
\author{}
\date{\today}

\begin{document}

\maketitle

\noindent\textbf{Hinweise:} Gesamtpunktzahl: \textbf{40 Punkte}. Die Aufgaben sind nach
Schwierigkeit aufsteigend sortiert: Verständnisfragen (Aufgaben 1--3), Standardaufgaben
(Aufgaben 4--6), Transfer-/Prüfungsniveau (Aufgaben 7--8). Erst selbst versuchen, Dinge
ausprobieren, damit spielen -- das ist kein Zeitrennen (vgl.\ Skript, S.~14). Taschenrechner
ist erlaubt; runden Sie, wo nicht anders angegeben, auf vier Nachkommastellen.

\section*{Aufgaben}

% ------------------------------------------------------------
\aufgabe{Analytisch vs.\ numerisch}{4 Punkte}\label{auf:1}
\begin{enumerate}[label=\alph*),itemsep=2pt]
  \item Definieren Sie in eigenen Worten, was numerische Methoden (numerical methods)
        sind, und nennen Sie zwei Gründe, warum man Lösungen approximativ statt analytisch
        exakt berechnet.
  \item Erläutern Sie kurz, warum numerische Methoden im Machine Learning -- insbesondere
        im Deep Learning -- unverzichtbar sind. Nennen Sie mindestens zwei Aspekte.
  \item Die Sigmoid-Funktion (sigmoid function) ist gegeben durch
        $\sigma(x) = \dfrac{1}{1 + e^{-x}}$.
        Erklären Sie, was eine transzendente Funktion (transcendental function) ist, und
        begründen Sie, warum die Gleichung $\sigma(x) = 0$ keine geschlossene Lösung
        (closed-form solution) besitzt. Was ist hier zusätzlich die Besonderheit der
        \glqq Nullstelle\grqq{} von $\sigma$?
\end{enumerate}

% ------------------------------------------------------------
\schreibraum[7]
\aufgabe{Diskretisierung und Schrittweite}{4 Punkte}\label{auf:2}
Eine kontinuierliche Funktion $f(x)$, $x \in [a,b]$, wird diskretisiert zu
$f(x_i)$ mit $x_i = a + ih$, $i = 0, 1, \ldots, N$, und Schrittweite
$h = \frac{b-a}{N}$.
\begin{enumerate}[label=\alph*),itemsep=2pt]
  \item Das Intervall $[0, 2]$ wird mit $N = 8$ Schritten diskretisiert. Berechnen Sie die
        Schrittweite $h$, geben Sie alle Gitterpunkte (grid points) $x_i$ an und bestimmen
        Sie deren Anzahl.
  \item Das Intervall $[-1, 1]$ soll mit Schrittweite $h = 0.01$ diskretisiert werden.
        Wie viele Schritte $N$ und wie viele Funktionsauswertungen sind nötig? Was passiert
        mit dem Rechenaufwand, wenn man $h$ halbiert?
  \item Warum ist die Wahl der Schrittweite $h$ wichtig? Beschreiben Sie den Zielkonflikt
        (trade-off) zwischen Genauigkeit und Rechenzeit, und nennen Sie einen
        Hauptunterschied zwischen einer kontinuierlichen Funktion und ihrer diskretisierten
        Version samt einer Implikation für das numerische Lösen.
\end{enumerate}

% ------------------------------------------------------------
\schreibraum[7]
\aufgabe{Extrema auf geschlossenen Intervallen}{3 Punkte}\label{auf:3}
\begin{enumerate}[label=\alph*),itemsep=2pt]
  \item Erklären Sie, warum bei der Suche nach Extrema einer Funktion auf einem
        geschlossenen Intervall (closed interval) $[a,b]$ neben den kritischen Punkten
        stets auch die Endpunkte $a$ und $b$ geprüft werden müssen.
  \item Bestimmen Sie das Minimum und das Maximum von $f(x) = x^2$ auf dem Intervall
        $[1, 3]$. Gehen Sie systematisch vor: kritische Punkte bestimmen, Lage prüfen,
        Kandidaten auswerten.
\end{enumerate}

% ------------------------------------------------------------
\schreibraum[6]
\aufgabe{Minimumsuche: analytisch und per Grid Search}{6 Punkte}\label{auf:4}
Gesucht ist das Minimum von $f(x) = \cos(x)$ auf dem geschlossenen Intervall $[0, 4]$.
\begin{enumerate}[label=\alph*),itemsep=2pt]
  \item Bestimmen Sie das Minimum \emph{analytisch}: Berechnen Sie die Ableitung, finden
        Sie alle kritischen Punkte im Intervall, klassifizieren Sie sie mit der zweiten
        Ableitung und vergleichen Sie mit den Werten an den Endpunkten.
  \item Bestimmen Sie das Minimum \emph{numerisch} per Brute-Force-Gittersuche
        (grid search): Diskretisieren Sie $[0,4]$ mit Schrittweite $h = 1$, werten Sie
        $f$ an allen Gitterpunkten aus (Tabelle!) und wählen Sie den Punkt mit dem
        kleinsten Funktionswert.
  \item Berechnen Sie den Approximationsfehler (approximation error) zwischen analytischem
        und numerischem Minimal\emph{wert}. Was würde mit dem Fehler geschehen, wenn man
        die Schrittweite verkleinert -- und was wäre der Preis dafür?
\end{enumerate}

% ------------------------------------------------------------
\schreibraum[9]
\aufgabe{Lineares $2\times 2$-System: Substitution vs.\ Grid Search}{6 Punkte}\label{auf:5}
Gegeben ist das lineare Gleichungssystem
\[
\begin{bmatrix} 3 & 1 \\ 4 & 2 \end{bmatrix}
\begin{bmatrix} w \\ b \end{bmatrix}
=
\begin{bmatrix} 9 \\ 10 \end{bmatrix}.
\]
\begin{enumerate}[label=\alph*),itemsep=2pt]
  \item Lösen Sie das System \emph{analytisch} per Substitution. Verifizieren Sie Ihre
        Lösung durch Einsetzen in beide Gleichungen.
  \item Lösen Sie das System \emph{numerisch} per Grid Search über das Gitter
        $w, b \in \{0,\ 2,\ 4\}$ (Schrittweite $2$). Verwenden Sie als Fehlermaß die Summe
        der absoluten Abweichungen
        \[
        \text{error}_1 = |3w + b - 9|, \qquad \text{error}_2 = |4w + 2b - 10|,
        \]
        werten Sie alle $3 \times 3 = 9$ Kombinationen tabellarisch aus und geben Sie das
        Paar $(w, b)$ mit dem kleinsten akkumulierten Fehler an.
  \item Vergleichen Sie das Grid-Search-Ergebnis mit der exakten Lösung aus a). Erklären
        Sie die Abweichung und gehen Sie dabei insbesondere darauf ein, welche Rolle die
        Wahl des Suchintervalls (vgl.\ Intervallwahl $[a,b]$) hier spielt.
\end{enumerate}

% ------------------------------------------------------------
\schreibraum[9]
\aufgabe{Lineares $3\times 3$-System analytisch lösen}{6 Punkte}\label{auf:6}
Lösen Sie das folgende System aus drei Gleichungen mit drei Unbekannten (Übungsaufgabe
aus dem Skript, Gl.~2, S.~10) per Substitution bzw.\ Elimination:
\[
\begin{bmatrix} 2 & 1 & 3 \\ 1 & 4 & 2 \\ 3 & 2 & 5 \end{bmatrix}
\begin{bmatrix} w_1 \\ w_2 \\ w_3 \end{bmatrix}
=
\begin{bmatrix} 14 \\ 20 \\ 32 \end{bmatrix}
\]
Geben Sie die Lösung exakt (als Brüche) an und verifizieren Sie sie in allen drei
Gleichungen. Notieren Sie anschließend in ein bis zwei Sätzen, was diese Rechnung im
Vergleich zum $2\times 2$-System über die Skalierung analytischer Lösungen
(computational complexity) lehrt.

% ------------------------------------------------------------
\schreibraum[9]
\aufgabe{Komplexität der Grid Search: $\mathcal{O}(N^d)$}{5 Punkte}\label{auf:7}
Eine Grid Search über $d$ Parameter mit je $N$ Gitterpunkten pro Parameter muss
$N^d$ Kombinationen auswerten (Ausblick: Komplexität $\mathcal{O}(N^d)$).
\begin{enumerate}[label=\alph*),itemsep=2pt]
  \item Im Skript wird ein $2\times 2$-System per Grid Search mit $5$ Gitterpunkten pro
        Variable gelöst. Wie viele Kombinationen sind das? Wie viele wären es bei
        $N = 100$ Gitterpunkten pro Variable und weiterhin $d = 2$?
  \item Ein kleines neuronales Netz habe $d = 10$ Parameter, jeder werde mit $N = 100$
        Gitterpunkten abgesucht. Wie viele Kombinationen entstehen? Wie lange dauert die
        Suche, wenn ein Computer $10^9$ Kombinationen pro Sekunde auswertet? Geben Sie das
        Ergebnis in Sekunden und in Jahren an (1 Jahr $\approx 3.156 \cdot 10^7$\,s).
  \item Folgern Sie: Warum ist Brute-Force-Grid-Search für das Training moderner
        Deep-Learning-Modelle (Millionen Parameter) ungeeignet, obwohl sie für das
        $2\times 2$-System gut funktioniert hat? Schlagen Sie -- als Denkanstoß im Sinne
        des Skript-Teasers (S.~16) -- eine einfache Idee vor, das Verhältnis von
        Genauigkeit zu Rechenaufwand der Grid Search zu verbessern.
\end{enumerate}

% ------------------------------------------------------------
\schreibraum[8]
\aufgabe{Transfer: Loss-Minimierung eines linearen Modells}{6 Punkte}\label{auf:8}
Ein lineares Modell $\hat{y} = xw$ (lineares Modell ohne Bias; die einfachste Form eines
neuronalen Netzes mit einer einzigen Einheit, vgl.\ Skript S.~14) soll an die Datenpunkte
\[
(x_1, y_1) = (1, 2), \qquad (x_2, y_2) = (2, 4), \qquad (x_3, y_3) = (3, 5)
\]
angepasst werden. Als Verlustfunktion (loss function) dient der quadratische Fehler
\[
L(w) = \sum_{i=1}^{3} (x_i w - y_i)^2 .
\]
\begin{enumerate}[label=\alph*),itemsep=2pt]
  \item Erklären Sie kurz, warum die Suche nach dem besten Parameter $w$ ein
        Minimierungsproblem ist und was das mit dem Training von ML-Modellen
        (Loss-Minimierung) zu tun hat.
  \item Führen Sie eine Grid Search über $w \in [0, 3]$ mit Schrittweite $h = 0.5$ durch:
        Werten Sie $L(w)$ an allen sieben Gitterpunkten aus (Tabelle) und geben Sie das
        numerisch gefundene Minimum an.
  \item Bestimmen Sie das Minimum von $L(w)$ auf $[0,3]$ nun \emph{analytisch}
        (Ableitung null setzen, zweite Ableitung prüfen, Endpunkte des Intervalls nicht
        vergessen). Geben Sie $w^\ast$ exakt und gerundet an, berechnen Sie $L(w^\ast)$
        und den Approximationsfehler $|w_{\text{grid}} - w^\ast|$ der Grid Search.
\end{enumerate}

\bigskip
\noindent\textbf{Summe: 40 Punkte}

% ============================================================
\schreibraum[9]
\newpage
\section*{Lösungen --- erst nach Bearbeitung lesen!}

% ------------------------------------------------------------
\paragraph{Lösung zu Aufgabe~\ref{auf:1}}\mbox{}\\[2pt]
\textbf{a)} Numerische Methoden sind ein Teilgebiet der Mathematik, in dem Lösungen nicht
analytisch exakt, sondern \emph{approximativ} berechnet werden. Zwei Gründe (Skript, S.~9):
(1)~Viele Probleme sind analytisch gar nicht lösbar oder zu komplex, um praktikabel zu sein;
(2)~man kann Rechenleistung (computational power) anzapfen, um approximative Lösungen
effizient zu erhalten.

\textbf{b)} In ML/AI sind numerische Methoden entscheidend für das Training von Modellen,
die Optimierung von Parametern und die Simulation komplexer Systeme, wo analytische
Lösungen unmöglich (infeasible) sind. Sie ermöglichen den effizienten Umgang mit großen
Datensätzen und komplexen Algorithmen, sodass Modelle effektiv aus Daten lernen können,
während die Rechenressourcen im Rahmen bleiben. Besonders im Deep Learning (Millionen
Parameter, umfangreiche Berechnungen) erleichtern sie die Optimierungsprozesse des
Modelltrainings (Skript, S.~9). Zwei beliebige dieser Aspekte genügen.

\textbf{c)} Eine transzendente Funktion kann nicht als endliche Kombination algebraischer
Operationen (Addition, Subtraktion, Multiplikation, Division, Wurzeln) ausgedrückt werden
und besitzt daher keine geschlossenen Lösungen (Skript, S.~11). Bei
$\sigma(x) = \frac{1}{1+e^{-x}}$ macht der Exponentialterm $e^{-x}$ es unmöglich, $x$ mit
elementaren Funktionen (Polynome, rationale, trigonometrische Funktionen) zu isolieren.
Besonderheit: Die Nullstelle existiert gar nicht -- $\sigma(x)$ nähert sich $0$ nur
asymptotisch für $x \to -\infty$, erreicht $0$ aber für keinen endlichen Wert von $x$
(Skript, Randnotiz S.~11). Es gibt also nicht nur keine geschlossene Lösung, sondern
überhaupt keine Lösung von $\sigma(x) = 0$.

% ------------------------------------------------------------
\paragraph{Lösung zu Aufgabe~\ref{auf:2}}\mbox{}\\[2pt]
\textbf{a)} Schrittweite nach Gl.~(4) des Skripts:
\[
h = \frac{b-a}{N} = \frac{2-0}{8} = 0.25 .
\]
Gitterpunkte $x_i = 0 + i \cdot 0.25$ für $i = 0, 1, \ldots, 8$:
\[
x_i \in \{0,\ 0.25,\ 0.5,\ 0.75,\ 1,\ 1.25,\ 1.5,\ 1.75,\ 2\}.
\]
Das sind $N + 1 = 9$ Gitterpunkte (Achtung: $N$ zählt die \emph{Schritte}, nicht die
Punkte).

\textbf{b)} Aus $h = \frac{b-a}{N}$ folgt
\[
N = \frac{b-a}{h} = \frac{1-(-1)}{0.01} = \frac{2}{0.01} = 200
\]
Schritte, also $N + 1 = 201$ Funktionsauswertungen. Halbiert man $h$ (auf $0.005$),
verdoppelt sich $N$ auf $400$ ($401$ Auswertungen) -- der Rechenaufwand wächst in einer
Dimension also (näherungsweise) linear mit $1/h$.

\textbf{c)} Kleines $h$ bedeutet ein feineres Gitter: höhere Genauigkeit der Approximation,
aber mehr Funktionsauswertungen und damit mehr Rechenzeit; großes $h$ ist billig, aber
ungenau. Die Wahl von $h$ steuert also direkt den Kompromiss zwischen Genauigkeit und
Rechenaufwand. Hauptunterschied kontinuierlich vs.\ diskretisiert: Die kontinuierliche
Funktion ist auf dem gesamten Intervall definiert, die diskretisierte Version nur an den
endlich vielen Gitterpunkten $x_i$ -- zwischen den Punkten hat man keine Information,
es gilt nur $f(x) \approx f(x_i)$ (Gl.~7 im Skript, S.~16). Implikation: Merkmale zwischen
den Gitterpunkten (z.\,B.\ das wahre Minimum) können verfehlt werden; jede numerische
Lösung trägt einen Approximationsfehler.

% ------------------------------------------------------------
\paragraph{Lösung zu Aufgabe~\ref{auf:3}}\mbox{}\\[2pt]
\textbf{a)} Auf einem geschlossenen Intervall $[a,b]$ kann ein Extremum an einem kritischen
Punkt (Ableitung null oder nicht definiert) im Inneren \emph{oder} an den Endpunkten $a$,
$b$ auftreten -- an den Endpunkten auch dann, wenn die Ableitung dort nicht null ist, weil
die Funktion dort schlicht \glqq abgeschnitten\grqq{} wird und kein innerer Vergleichspunkt
mehr existiert. Wer nur kritische Punkte prüft, kann das tatsächliche Extremum übersehen
(Skript, S.~13).

\textbf{b)} $f(x) = x^2$, $f'(x) = 2x$. Kritischer Punkt: $2x = 0 \Rightarrow x = 0$.
Aber $0 \notin [1,3]$ -- es gibt also \emph{keinen} kritischen Punkt im Intervall. Die
Extrema liegen daher zwingend an den Endpunkten:
\[
f(1) = 1, \qquad f(3) = 9 .
\]
Minimum: $f(1) = 1$ bei $x = 1$; Maximum: $f(3) = 9$ bei $x = 3$ -- obwohl die Ableitung
an beiden Stellen ($f'(1) = 2 \neq 0$, $f'(3) = 6 \neq 0$) nicht verschwindet. Genau das
ist die Stolperfalle aus Teil a).

% ------------------------------------------------------------
\paragraph{Lösung zu Aufgabe~\ref{auf:4}}\mbox{}\\[2pt]
\textbf{a) Analytisch.} $f(x) = \cos(x)$, $f'(x) = -\sin(x)$. Kritische Punkte:
\[
-\sin(x) = 0 \iff x = k\pi, \quad k \in \mathbb{Z}.
\]
Im Intervall $[0,4]$ liegen $x = 0$ und $x = \pi \approx 3.1416$ (der nächste Kandidat
$2\pi \approx 6.2832 > 4$ liegt außerhalb). Zweite Ableitung: $f''(x) = -\cos(x)$.
\[
f''(0) = -\cos(0) = -1 < 0 \ \text{(lokales Maximum)}, \qquad
f''(\pi) = -\cos(\pi) = 1 > 0 \ \text{(lokales Minimum)}.
\]
Kandidaten auswerten (kritische Punkte \emph{und} Endpunkte):
\[
f(0) = \cos(0) = 1, \qquad f(\pi) = \cos(\pi) = -1, \qquad f(4) = \cos(4) \approx -0.6536 .
\]
\textbf{Ergebnis:} Das Minimum ist $-1$ bei $x = \pi \approx 3.1416$.

\textbf{b) Numerisch (Grid Search, $h = 1$).} Mit $h = \frac{4-0}{N} = 1$ ist $N = 4$,
also fünf Gitterpunkte $x_i = 0 + i \cdot 1$:
\begin{center}
\begin{tabular}{ll}
\toprule
$x_i$ & $f(x_i) = \cos(x_i)$ \\
\midrule
$x_0 = 0$ & $\cos(0) = 1$ \\
$x_1 = 1$ & $\cos(1) \approx 0.5403$ \\
$x_2 = 2$ & $\cos(2) \approx -0.4161$ \\
$x_3 = 3$ & $\cos(3) \approx -0.9900$ \\
$x_4 = 4$ & $\cos(4) \approx -0.6536$ \\
\bottomrule
\end{tabular}
\end{center}
Das Minimum unter diesen Werten ist $\cos(3) \approx -0.9900$ bei $x = 3$.

\textbf{c) Approximationsfehler.} Differenz zwischen analytischem und numerischem
Minimalwert:
\[
\left| -1 - (-0.9900) \right| = 0.0100 .
\]
Eine kleinere Schrittweite legt Gitterpunkte näher an das wahre Minimum bei $x = \pi$ und
verkleinert den Fehler (z.\,B.\ träfe $h = 0.5$ den Punkt $x = 3$ weiterhin, $h = 0.1$
schon $x = 3.1$ mit $\cos(3.1) \approx -0.9991$). Der Preis: mehr Gitterpunkte, also mehr
Funktionsauswertungen und mehr Rechenzeit -- der bekannte Genauigkeit-Rechenaufwand-Kompromiss.

% ------------------------------------------------------------
\paragraph{Lösung zu Aufgabe~\ref{auf:5}}\mbox{}\\[2pt]
\textbf{a) Substitution.} Die Gleichungen lauten
\[
\text{(I)}\ 3w + b = 9, \qquad \text{(II)}\ 4w + 2b = 10 .
\]
Aus (I): $b = 9 - 3w$. Einsetzen in (II):
\[
4w + 2(9 - 3w) = 10
\;\Rightarrow\;
4w + 18 - 6w = 10
\;\Rightarrow\;
-2w = -8
\;\Rightarrow\;
w = 4 .
\]
Rückeinsetzen: $b = 9 - 3 \cdot 4 = 9 - 12 = -3$.
\textbf{Lösung:} $(w, b) = (4, -3)$.
Verifikation: (I): $3 \cdot 4 + (-3) = 12 - 3 = 9$ \checkmark;\quad
(II): $4 \cdot 4 + 2 \cdot (-3) = 16 - 6 = 10$ \checkmark.

\textbf{b) Grid Search.} Für jede Kombination $(w,b)$ aus $\{0, 2, 4\}^2$ werden
$\text{error}_1 = |3w + b - 9|$ und $\text{error}_2 = |4w + 2b - 10|$ sowie der
akkumulierte Fehler (Summe) berechnet:
\begin{center}
\begin{tabular}{rrllr}
\toprule
$w$ & $b$ & $3w+b$ (error$_1$) & $4w+2b$ (error$_2$) & Error \\
\midrule
0 & 0 & 0 (9)   & 0 (10)  & 19 \\
0 & 2 & 2 (7)   & 4 (6)   & 13 \\
0 & 4 & 4 (5)   & 8 (2)   & 7 \\
2 & 0 & 6 (3)   & 8 (2)   & 5 \\
2 & 2 & 8 (1)   & 12 (2)  & \textbf{3*} \\
2 & 4 & 10 (1)  & 16 (6)  & 7 \\
4 & 0 & 12 (3)  & 16 (6)  & 9 \\
4 & 2 & 14 (5)  & 20 (10) & 15 \\
4 & 4 & 16 (7)  & 24 (14) & 21 \\
\bottomrule
\end{tabular}
\end{center}
Der minimale akkumulierte Fehler ist $3$ bei $(w, b) = (2, 2)$ (mit * markiert).

\textbf{c) Vergleich und Diskussion.} Die exakte Lösung $(4, -3)$ wird von der Grid Search
deutlich verfehlt: Das Gitter liefert $(2, 2)$. Zwei Gründe: Erstens ist das Gitter mit
Schrittweite $2$ sehr grob. Zweitens -- und entscheidend -- liegt die wahre Lösung in der
$b$-Richtung ($b = -3$) \emph{außerhalb} des Suchbereichs $[0, 4]$: Egal wie fein man
dieses Gitter macht, die Grid Search kann die exakte Lösung nie finden, weil sie nur
innerhalb des gewählten Intervalls sucht. Das illustriert die Selbstreflexionsfrage des
Skripts (S.~15): Die Wahl des Intervalls $[a,b]$ beeinflusst, welche (Pseudo-)Optima
numerisch überhaupt gefunden werden können -- die beste Gitterlösung ist nur die beste
\emph{innerhalb des Gitters}.

% ------------------------------------------------------------
\paragraph{Lösung zu Aufgabe~\ref{auf:6}}\mbox{}\\[2pt]
Die Gleichungen lauten
\[
\text{(I)}\ 2w_1 + w_2 + 3w_3 = 14, \qquad
\text{(II)}\ w_1 + 4w_2 + 2w_3 = 20, \qquad
\text{(III)}\ 3w_1 + 2w_2 + 5w_3 = 32 .
\]
\textbf{Schritt 1:} (I) nach $w_2$ auflösen:
\[
w_2 = 14 - 2w_1 - 3w_3 .
\]
\textbf{Schritt 2:} In (II) einsetzen:
\[
w_1 + 4(14 - 2w_1 - 3w_3) + 2w_3 = 20
\;\Rightarrow\;
w_1 + 56 - 8w_1 - 12w_3 + 2w_3 = 20
\;\Rightarrow\;
-7w_1 - 10w_3 = -36,
\]
also $7w_1 + 10w_3 = 36$ \ (II$'$).

\textbf{Schritt 3:} In (III) einsetzen:
\[
3w_1 + 2(14 - 2w_1 - 3w_3) + 5w_3 = 32
\;\Rightarrow\;
3w_1 + 28 - 4w_1 - 6w_3 + 5w_3 = 32
\;\Rightarrow\;
-w_1 - w_3 = 4,
\]
also $w_1 = -4 - w_3$ \ (III$'$).

\textbf{Schritt 4:} (III$'$) in (II$'$) einsetzen:
\[
7(-4 - w_3) + 10w_3 = 36
\;\Rightarrow\;
-28 - 7w_3 + 10w_3 = 36
\;\Rightarrow\;
3w_3 = 64
\;\Rightarrow\;
w_3 = \frac{64}{3} .
\]
\textbf{Schritt 5:} Rückwärts einsetzen:
\[
w_1 = -4 - \frac{64}{3} = \frac{-12 - 64}{3} = -\frac{76}{3}, \qquad
w_2 = 14 - 2\left(-\frac{76}{3}\right) - 3 \cdot \frac{64}{3}
    = 14 + \frac{152}{3} - 64 = \frac{42 + 152 - 192}{3} = \frac{2}{3} .
\]
\textbf{Lösung:}
\[
w_1 = -\frac{76}{3} \approx -25.3333, \qquad
w_2 = \frac{2}{3} \approx 0.6667, \qquad
w_3 = \frac{64}{3} \approx 21.3333 .
\]
\textbf{Verifikation:}
\[
\text{(I)}:\ 2\left(-\tfrac{76}{3}\right) + \tfrac{2}{3} + 3 \cdot \tfrac{64}{3}
= \tfrac{-152 + 2 + 192}{3} = \tfrac{42}{3} = 14\ \checkmark
\]
\[
\text{(II)}:\ -\tfrac{76}{3} + 4 \cdot \tfrac{2}{3} + 2 \cdot \tfrac{64}{3}
= \tfrac{-76 + 8 + 128}{3} = \tfrac{60}{3} = 20\ \checkmark
\]
\[
\text{(III)}:\ 3\left(-\tfrac{76}{3}\right) + 2 \cdot \tfrac{2}{3} + 5 \cdot \tfrac{64}{3}
= \tfrac{-228 + 4 + 320}{3} = \tfrac{96}{3} = 32\ \checkmark
\]
\textbf{Lehre zur Skalierung:} Schon der Schritt von $2$ auf $3$ Unbekannte verlangt
deutlich mehr Substitutionsschritte, Zwischenterme und Sorgfalt (Bruchrechnung!). Mit
wachsender Systemgröße (Tausende Gleichungen mit Tausenden Unbekannten) werden analytische
Handlösungen wegen der Rechenkomplexität und Zeitbeschränkungen unpraktikabel -- genau das
Argument des Skripts (S.~10--11) für numerische, computergestützte Methoden.

% ------------------------------------------------------------
\paragraph{Lösung zu Aufgabe~\ref{auf:7}}\mbox{}\\[2pt]
\textbf{a)} Mit $N = 5$ Gitterpunkten und $d = 2$ Variablen:
$N^d = 5^2 = 25$ Kombinationen (genau die 25 Zeilen der Tabelle~1 im Skript, S.~15; das
Skript spricht von \glqq only 100 operations\grqq, also etwa vier Rechenoperationen pro
Kombination). Mit $N = 100$ und $d = 2$: $100^2 = 10\,000$ Kombinationen -- für einen
Computer immer noch trivial.

\textbf{b)} Mit $N = 100$ und $d = 10$:
\[
N^d = 100^{10} = (10^2)^{10} = 10^{20} \ \text{Kombinationen}.
\]
Bei $10^9$ Auswertungen pro Sekunde:
\[
t = \frac{10^{20}}{10^9} = 10^{11}\ \text{s}
\qquad\Rightarrow\qquad
\frac{10^{11}}{3.156 \cdot 10^7\ \text{s/Jahr}} \approx 3.17 \cdot 10^3\ \text{Jahre},
\]
also rund $3\,170$ Jahre -- für ein \glqq kleines\grqq{} Netz mit nur zehn Parametern.

\textbf{c)} Die Anzahl der Kombinationen wächst \emph{exponentiell} in der Zahl der
Parameter $d$ (Fluch der Dimensionalität, curse of dimensionality). Deep-Learning-Modelle
haben Millionen Parameter; $N^{d}$ ist dann astronomisch groß und Brute Force vollkommen
aussichtslos -- deshalb braucht das Modelltraining effizientere numerische
Optimierungsverfahren (Ausblick des Kurses, z.\,B.\ ableitungsbasierte Methoden).
Idee zum Teaser (S.~16): \emph{adaptive bzw.\ verfeinerte Suche} -- erst mit grobem Gitter
suchen, dann nur in der Umgebung des besten Kandidaten mit feinerem Gitter weitersuchen.
So erreicht man die Genauigkeit eines feinen Gitters zu einem Bruchteil der Auswertungen
(in 1D z.\,B.\ $10 + 10 = 20$ statt $100$ Auswertungen für die gleiche Endauflösung).

% ------------------------------------------------------------
\paragraph{Lösung zu Aufgabe~\ref{auf:8}}\mbox{}\\[2pt]
\textbf{a)} Der Loss $L(w)$ misst, wie stark die Modellvorhersagen $\hat{y}_i = x_i w$ von
den beobachteten Werten $y_i$ abweichen. Der beste Parameter ist der mit der kleinsten
Gesamtabweichung -- die Suche nach $w$ ist also die Minimierung von $L(w)$. Genau das ist
Training im Machine Learning: Modellparameter werden so eingestellt, dass eine
Verlustfunktion minimal wird; $xw$ (bzw.\ $xw + b$) ist dabei die einfachste Form eines
neuronalen Netzes mit einer einzigen Einheit (Skript, S.~14). Die Grid Search aus b) ist
die primitivste numerische Methode dafür.

\textbf{b) Grid Search} über $w \in \{0, 0.5, 1, 1.5, 2, 2.5, 3\}$ mit
$L(w) = (w - 2)^2 + (2w - 4)^2 + (3w - 5)^2$:
\begin{center}
\begin{tabular}{lllll}
\toprule
$w$ & $(w-2)^2$ & $(2w-4)^2$ & $(3w-5)^2$ & $L(w)$ \\
\midrule
$0$   & $4$    & $16$  & $25$    & $45$ \\
$0.5$ & $2.25$ & $9$   & $12.25$ & $23.5$ \\
$1$   & $1$    & $4$   & $4$     & $9$ \\
$1.5$ & $0.25$ & $1$   & $0.25$  & $1.5$ \\
$2$   & $0$    & $0$   & $1$     & $\mathbf{1\ast}$ \\
$2.5$ & $0.25$ & $1$   & $6.25$  & $7.5$ \\
$3$   & $1$    & $4$   & $16$    & $21$ \\
\bottomrule
\end{tabular}
\end{center}
Numerisch gefundenes Minimum: $w_{\text{grid}} = 2$ mit $L(2) = 1$ (mit * markiert).

\textbf{c) Analytisch.} Ausdifferenzieren (Kettenregel):
\[
L'(w) = 2(w - 2) \cdot 1 + 2(2w - 4) \cdot 2 + 2(3w - 5) \cdot 3
      = 2w - 4 + 8w - 16 + 18w - 30 = 28w - 50 .
\]
Nullsetzen:
\[
28w - 50 = 0 \;\Rightarrow\; w^\ast = \frac{50}{28} = \frac{25}{14} \approx 1.7857 .
\]
Zweite Ableitung: $L''(w) = 28 > 0$, also Minimum; $w^\ast \in [0,3]$ liegt im Intervall.
Endpunkte prüfen: $L(0) = 45$ und $L(3) = 21$ sind beide größer -- das Minimum auf $[0,3]$
liegt also tatsächlich bei $w^\ast$. Minimaler Loss:
\[
L\!\left(\tfrac{25}{14}\right)
= \left(\tfrac{25}{14} - \tfrac{28}{14}\right)^2
+ \left(\tfrac{50}{14} - \tfrac{56}{14}\right)^2
+ \left(\tfrac{75}{14} - \tfrac{70}{14}\right)^2
= \tfrac{9}{196} + \tfrac{36}{196} + \tfrac{25}{196}
= \tfrac{70}{196} = \tfrac{5}{14} \approx 0.3571 .
\]
Approximationsfehler der Grid Search im Parameter:
\[
|w_{\text{grid}} - w^\ast| = \left|2 - \tfrac{25}{14}\right| = \tfrac{3}{14} \approx 0.2143 ,
\]
und im Loss: $|1 - \tfrac{5}{14}| = \tfrac{9}{14} \approx 0.6429$. Ein feineres Gitter
(kleineres $h$) würde beide Fehler verkleinern -- zum Preis von mehr Auswertungen.

\end{document}
